% The official checklist ships a \item[] Guidelines: itemize under every question. Those are
% instructions ADDRESSED TO THE AUTHOR, not content: 2,486 of this file's 4,185 words (59%).
% Stripped, because this is a WORKSHOP submission where the checklist is not required at all
% (main-track requires it; workshops normally waive it and this CFP is silent), and 6 physical
% pages of self-addressed boilerplate is not what the appendix is for. Every QUESTION and every
% ANSWER/JUSTIFICATION is untouched.
% ==> IF THIS PAPER GOES TO AN ICML/NeurIPS MAIN TRACK, restore them from git history: they are
%     part of the official template there and the answers below are already written.
\section*{NeurIPS Paper Checklist}




\begin{enumerate}

\item {\bf Claims}
    \item[] Question: Do the main claims made in the abstract and introduction accurately reflect the paper's contributions and scope?
    \item[] Answer: \answerYes{}
    \item[] Justification: The abstract and \S\ref{sec:ident} state both the corrective claim (ally-fit ``truth'' probes are unidentified) and the constructive one (\tm{} returns to $1.000$), and both are scoped in place: the headline conjunction of a saturated policy \emph{and} a wholly unmodified recipe is stated as one family (\S\ref{sec:constructive}), the \tm{} result as 19 of 21 identification cells with both exceptions named (the separate set of 21 emergent grid cells that carries the saturation nulls is a different denominator and is labelled as such), and the causal claim scoped to the axis and the outcome actually tested: the criticized ally-fit direction is the \emph{potent} lever and the identified mixed-fit direction largely inert (App.~\ref{app:causal}, \S\ref{sec:signed}), established by a paired McNemar test on matched episodes rather than offered as suggestive. The scope of that claim is answer choice, not truth: on the two arms where neither rate sits on a floor the same push raises the true-bit-$1$ answer rate and \emph{lowers} the true-bit-$0$ one, which is what the paper concludes it is, one token-preference knob rather than a truth axis, so no truth-steering capability is claimed anywhere.
\item {\bf Limitations}
    \item[] Question: Does the paper discuss the limitations of the work performed by the authors?
    \item[] Answer: \answerYes{}
    \item[] Justification: \S\ref{sec:limits} is a dedicated section covering the single one-token game, the one-of-four-families saturation scope, the floored subpopulations, single-layer and single-site steering, the settling experiment's imperfect and strongly depth-dependent transfer, and the crossed text-versus-behaviour design that this probe site cannot support. Negative and non-replicating results are \emph{summarised in the body}, inside that same section, with the per-result detail in App.~\ref{app:negatives}: the body names them, so a reader who stops at page nine still sees that two of our own hypotheses were refuted, and the appendix carries the evidence rather than the disclosure.
\item {\bf Theory assumptions and proofs}
    \item[] Question: For each theoretical result, does the paper provide the full set of assumptions and a complete (and correct) proof?
    \item[] Answer: \answerYes{}
    \item[] Justification: The two formal results are stated with their assumptions and derived in place rather than in numbered theorem environments: the forced identity $\actally = 1 - \ta$ follows from the ally-data label vectors for \texttt{truth} and \texttt{action} being identical plus the antisymmetry of AUROC under label negation (\S\ref{sec:ident}), and is additionally verified numerically at all 751 (cell, layer) pairs to a maximum deviation of $2.2\times10^{-16}$; the aliasing audit $\gamma$ (Eq.~\ref{eq:gamma}) is derived in one line in App.~\ref{app:instrpair} under the stated honest-rung assumption.
    \item {\bf Experimental result reproducibility}
    \item[] Question: Does the paper fully disclose all the information needed to reproduce the main experimental results of the paper to the extent that it affects the main claims and/or conclusions of the paper (regardless of whether the code and data are provided or not)?
    \item[] Answer: \answerYes{}
    \item[] Justification: \S\ref{sec:setup} gives the game and reward structure (REINFORCE; ally $+1/-1$, rival $+1.5/-2$), the adapter configuration ($r{=}16$, $\alpha{=}32$, on q/k/v/o), the exact base-model versions, dtype, the \texttt{transformers} version, $N$ per cell, the read position, the three label sets, the two fitting regimes and the held-out ally / held-out rival scoring split. Where a number depends on a choice we made (chance tolerance, permutation count, $\alpha$ normalisation) the choice is stated with the number.
\item {\bf Open access to data and code}
    \item[] Question: Does the paper provide open access to the data and code, with sufficient instructions to faithfully reproduce the main experimental results, as described in supplemental material?
    \item[] Answer: \answerYes{}
    \item[] Justification: We release the game, the randomized-codebook variant, the RL recipe, all probe-fitting regimes, the intervention code and the figure pipeline, together with the exact per-run commands for every cell reported here (\S\ref{sec:repro}). For review the code is provided anonymised at the repository link given in \S\ref{sec:repro}. The figures in this paper are rendered from a consolidated results file that ships with the code, so every plotted number is re-derivable without re-running a GPU job. We deliberately do \emph{not} claim a one-command hosted reproduction path: reproducing the underlying runs requires an accelerator of the class named in \S\ref{sec:setup}, and we state the commands rather than ship a wrapper we have not executed end to end. Independently of the release, the paper states the protocol in enough detail to reproduce the main results from scratch (question 4).
\item {\bf Experimental setting/details}
    \item[] Question: Does the paper specify all the training and test details (e.g., data splits, hyperparameters, how they were chosen, type of optimizer) necessary to understand the results?
    \item[] Answer: \answerYes{}
    \item[] Justification: Measurement-side settings are fully specified in the paper. Probe type (per-layer logistic), read position (final prompt position, the answer-prediction position), layer grids, $N$ per cell, the ally-only versus mixed ally+rival fitting regimes and the held-out scoring split are in \S\ref{sec:setup}; the intervention settings (rank ladders, fit sizes relative to the residual width, dose grids in $\alpha_{\mathrm{rel}}$ rather than raw $\alpha$) are given where used in \S\ref{sec:signed} and \S\ref{sec:depth}. Since every claim in the paper is a claim about a \emph{probe} or an \emph{intervention}, these are the settings the reported numbers are a function of. Training-side settings are split between the paper and the release, and we state which is which. \S\ref{sec:setup} gives the RL algorithm (REINFORCE), the reward values, the adapter configuration ($r{=}16$, $\alpha{=}32$, on q/k/v/o) and a qualitative stopping rule (``until reward plateaus''); the gradient optimizer, learning rate, batch size and epoch budget are \emph{not} in the paper and are in the released per-run commands (question 5), which is where a replicator obtains them. One gap is visible in the paper and we flag it rather than let a reader assume the text is complete: Table~\ref{tab:ident} has a row named ``larger step size / batch'' whose two values appear nowhere in the text. That row is a robustness arm rather than a load-bearing cell, being one of five recipes that all give the same contrast, so the omission does not affect a reported conclusion.
\item {\bf Experiment statistical significance}
    \item[] Question: Does the paper report error bars suitably and correctly defined or other appropriate information about the statistical significance of the experiments?
    \item[] Answer: \answerYes{}
    \item[] Justification: Where intervals exist they are episode-level bootstraps; ablations are judged against \emph{permutation-calibrated} nulls rather than a fixed band alone, and cells clearing the tolerance but sitting above their permutation null are reported as suggestive, not established (\S\ref{sec:signed}). We also state the limits of our own error bars: the bootstrap holds the fitted probe and the partition fixed, so it is not a run-to-run interval, and we quantify refit variance separately across five split seeds in App.~\ref{app:objections} (a new held-out wording set and a new episode partition each time), where the across-refit sd is small at some depths and an order of magnitude larger at others, which is why the layer a claim is anchored at is stated rather than left implicit; \S\ref{sec:refit}, by contrast, reports a resample of the ally \emph{subset} and is a different quantity; one equivalence test is reported as \emph{undefined} because its paired difference has exactly zero variance, and single-run cells without intervals are named as such.
\item {\bf Experiments compute resources}
    \item[] Question: For each experiment, does the paper provide sufficient information on the computer resources (type of compute workers, memory, time of execution) needed to reproduce the experiments?
    \item[] Answer: \answerYes{}
    \item[] Justification: \S\ref{sec:setup} states the type of compute worker (a single H200-class GPU per run), dtype (bf16), the software version and the per-cell $N$, which fixes the scale of every reported cell and is what a replicator needs in order to know what hardware to obtain. Two of the three things the question names are not reported, and we state that rather than leave it inferred: we give no GPU memory footprint and no wall-clock time for any experiment, and no aggregate, neither total GPU-hours nor the cost of the preliminary and failed runs, of which there were many, so the figure a replicator would budget is larger than the sum of the cells in the paper and we cannot say by how much. The two omissions are of different weight. Memory is bounded above by the accelerator named, so the reported worker type already constrains it; time of execution was not instrumented, and we will not report a figure we did not measure.
\item {\bf Code of ethics}
    \item[] Question: Does the research conducted in the paper conform, in every respect, with the NeurIPS Code of Ethics \url{https://neurips.cc/public/EthicsGuidelines}?
    \item[] Answer: \answerYes{}
    \item[] Justification: We have reviewed the NeurIPS Code of Ethics and are not aware of any respect in which this work departs from it. No human subjects, no personal data, and no scraped content are involved.
\item {\bf Broader impacts}
    \item[] Question: Does the paper discuss both potential positive societal impacts and negative societal impacts of the work performed?
    \item[] Answer: \answerYes{}
    \item[] Justification: The paper's central result is itself the safety-relevant impact and is discussed as such: a probing protocol already proposed for deception monitoring can return confidently inverted readings that in-distribution validation cannot detect (\S\ref{sec:ident}, App.~\ref{app:instrpair}), so the negative impact of the \emph{status quo} is false confidence in a deployed monitor. The positive side is the constructive fix and five prescriptions. We scope the contribution to the identifiability of the measurement rather than to detector deployment (\S\ref{sec:limits}), which is the claim that could otherwise be over-read.
\item {\bf Safeguards}
    \item[] Question: Does the paper describe safeguards that have been put in place for responsible release of data or models that have a high risk for misuse (e.g., pre-trained language models, image generators, or scraped datasets)?
    \item[] Answer: \answerNA{}
    \item[] Justification: The artifacts pose no realistic misuse risk requiring safeguards: the deceptive policies are adapters trained on a synthetic one-bit guessing game with a single-token answer, and they confer no capability that transfers outside it. No trained artifact is released: no model or adapter weights, no fitted probe or steering vector, and no corpus of model-generated text (the results file that does ship, question 13, holds aggregate measurements such as AUROCs and rates, not generations), so nothing ships that could be repurposed as a deceptive policy or as an off-the-shelf detector. What does ship is \emph{code} (questions 5 and 13), including the probe-fitting and intervention code; running it requires supplying one's own model and one's own compute, and it produces a measurement rather than a capability.
\item {\bf Licenses for existing assets}
    \item[] Question: Are the creators or original owners of assets (e.g., code, data, models), used in the paper, properly credited and are the license and terms of use explicitly mentioned and properly respected?
    \item[] Answer: \answerYes{}
    \item[] Justification: Every base model is credited by exact name and version in \S\ref{sec:setup}, and App.~\ref{app:licenses} names the governing license for each one individually, noting that they are not uniform: Mistral-7B-v0.3 and Qwen2.5-\{7,14,32\}B are Apache-2.0, whereas Llama-3.1-8B is under the Llama 3.1 Community License, Gemma-2-9B under the Gemma Terms of Use, and Qwen2.5-3B under the Qwen Research License rather than Apache-2.0: Qwen exempts its 3B and 72B variants. All use here is non-commercial research, which every one of these licenses permits.
\item {\bf New assets}
    \item[] Question: Are new assets introduced in the paper well documented and is the documentation provided alongside the assets?
    \item[] Answer: \answerYes{}
    \item[] Justification: There are two new assets, both provided anonymised at the repository link in \S\ref{sec:repro} (question 5). \emph{(i) Code}: the game and its randomized-codebook variant, the RL recipe, the probe-fitting and intervention code and the figure pipeline, with the exact per-run command for every cell reported here. \emph{(ii) Data}: a single consolidated results file holding the aggregate measurements (AUROCs, deception rates, cosines, dose responses) that every figure and table in this paper is rendered from, so a reader can re-derive any plotted number without a GPU. Documentation accompanies both: a description of the package layout, the per-run commands, and, for the results file specifically, the figure pipeline that consumes it ships alongside it, so the mapping from a stored field to a printed number is executable rather than merely described. The terms under which both assets are released are declared in the package itself, where a user will look for them, and no third-party asset is redistributed inside it: no base-model files, no weights and no adapters, so ours are the only terms that need declaring. \textbf{No trained weights or adapters are released}, so the redistribution and naming obligations that App.~\ref{app:licenses} records for the Llama and Gemma terms do not attach to this release.
\item {\bf Crowdsourcing and research with human subjects}
    \item[] Question: For crowdsourcing experiments and research with human subjects, does the paper include the full text of instructions given to participants and screenshots, if applicable, as well as details about compensation (if any)? 
    \item[] Answer: \answerNA{}
    \item[] Justification: The work involves no crowdsourcing and no human subjects. All data are generated by the synthetic game described in \S\ref{sec:setup}.
\item {\bf Institutional review board (IRB) approvals or equivalent for research with human subjects}
    \item[] Question: Does the paper describe potential risks incurred by study participants, whether such risks were disclosed to the subjects, and whether Institutional Review Board (IRB) approvals (or an equivalent approval/review based on the requirements of your country or institution) were obtained?
    \item[] Answer: \answerNA{}
    \item[] Justification: No research with human subjects was conducted, so no IRB or equivalent review was required.
\item {\bf Declaration of LLM usage}
    \item[] Question: Does the paper describe the usage of LLMs if it is an important, original, or non-standard component of the core methods in this research? Note that if the LLM is used only for writing, editing, or formatting purposes and does \emph{not} impact the core methodology, scientific rigor, or originality of the research, declaration is not required.
    %this research? 
    \item[] Answer: \answerNA{}
    \item[] Justification: LLMs are the \emph{object of study} here, not a component of the method: the models named in \S\ref{sec:setup} are what we train and probe. No LLM plays an important, original or non-standard role in the core methodology, which is probe fitting, activation interventions and the statistical tests reported above.
\end{enumerate}